\documentclass[11pt,a4paper]{article}

% ---------------------------------------------------------
% PACKAGES
% ---------------------------------------------------------
\usepackage[utf8]{inputenc}
\usepackage[T1]{fontenc}
\usepackage{lmodern}
\usepackage{amsmath, amssymb, amsfonts, bm}
\usepackage{geometry}
\usepackage{graphicx}
\usepackage{hyperref}
\usepackage{authblk}
\usepackage{cite}
\usepackage{physics}
\usepackage{mathtools}
\usepackage{tensor}
\usepackage{enumitem}

\geometry{margin=2.4cm}

% ---------------------------------------------------------
% TITLE
% ---------------------------------------------------------
\title{\textbf{Companion II (v49.1)}\\[0.2cm]
\textbf{Global Tensor Coherence in a CPT-Symmetric Bipartite Universe}\\[0.2cm]
\large Foundations of Relaxation-Driven Cyclic Cosmology (RDCC)}

\author[1]{Michael Lehmann}
\affil[1]{\small Independent Researcher, Relaxation-Driven Cyclic Cosmology (RDCC)\\
\texttt{mi.lehmann@gmx.de}}

\date{September 2026}

\begin{document}
\maketitle

% ========================================================
% ABSTRACT
% ========================================================

\begin{abstract}
This Companion develops the tensor sector of RDCC in its fully reconstructed
and Phase-1-integrated v49.1 form. Tensor modes are global geometric degrees of
freedom defined on the bipartite Hilbert space and remain coherent across both
sectors of the CPT-symmetric global state. Their evolution is governed not by
local field dynamics but by the global structure of the metric. This Companion
derives the corrected propagation equations, establishes the role of the
dimension-8 inter-sector operator, and explains the quadratic suppression and
logarithmic modulation of the stochastic gravitational-wave background. These
results integrate the relaxation sector (Companion III), the ultraviolet sector
(Companion X), and the tensor–relaxation coupling (Companion XI).
\end{abstract}

\tableofcontents
\newpage
% ========================================================
% 1 INTRODUCTION
% ========================================================

\section{Introduction}

In standard cosmology, tensor modes are treated as local perturbations of the
metric, generated during inflation and subsequently decohering as the Universe
expands. RDCC replaces this picture with a fundamentally different structure:
the Universe is described by a global CPT-symmetric quantum state defined on a
bipartite Hilbert space. The metric is a global geometric degree of freedom
that couples to both sectors simultaneously.

As a consequence, tensor modes are not local fields but global excitations of
the metric. They remain coherent across both sectors and are not subject to the
decoherence that affects scalar and vector modes. This global coherence is a
structural prediction of RDCC and forms the foundation for the tensor sector.

This Companion develops the corrected tensor sector of RDCC in its v49.1 form.
We derive the global tensor equation, establish the role of the dimension-8
inter-sector operator, and explain the origin of the quadratic suppression and
logarithmic modulation of the stochastic gravitational-wave background.
% ========================================================
% 2 GLOBAL TENSOR STRUCTURE
% ========================================================

\section{Global Tensor Structure}

Tensor modes are defined as transverse-traceless perturbations of the global
metric,
\begin{equation}
g_{\mu\nu} = \bar{g}_{\mu\nu} + h_{\mu\nu},
\end{equation}
with
\begin{equation}
\nabla^\mu h_{\mu\nu} = 0, \qquad h^\mu_{\ \mu} = 0.
\end{equation}

Unlike scalar and vector modes, tensor modes are not associated with local
degrees of freedom. They are global excitations of the metric and therefore
remain coherent across both sectors of the bipartite Hilbert space.
% --------------------------------------------------------
% 2.1 Bipartite Coherence
% --------------------------------------------------------

\subsection[Bipartite Tensor Coherence]{Bipartite Tensor Coherence}

The global state is written as
\begin{equation}
\ket{\Psi} = \sum_n \lambda_n \ket{n}_{(+)} \otimes \ket{n}_{(-)}.
\end{equation}

Tensor modes correspond to excitations of the global metric operator,
\begin{equation}
\hat{h}_{\mu\nu}
=
\hat{h}_{\mu\nu}^{(+)} \otimes \mathbb{1}_{(-)}
+
\mathbb{1}_{(+)} \otimes \hat{h}_{\mu\nu}^{(-)}.
\end{equation}

Because the metric is global, tensor modes do not decohere under sectoral
projection. The reduced density matrices
\begin{equation}
\rho_{(+)} = \Tr_{(-)} \ket{\Psi}\bra{\Psi}, \qquad
\rho_{(-)} = \Tr_{(+)} \ket{\Psi}\bra{\Psi},
\end{equation}
retain the full tensor coherence of the global state.
% ========================================================
% 3 DIMENSION-8 OPERATOR AND TENSOR MODULATION
% ========================================================

\section[Dimension-8 Operator and Tensor Modulation]{Dimension-8 Operator and Tensor Modulation}

The inter-sector operator
\begin{equation}
\mathcal{O}_{\mathrm{int}} = T^{(+)}_{\mu\nu} T_{(-)}^{\mu\nu}
\end{equation}
has mass dimension 8. Its contribution to the tensor sector is suppressed by
$M^{-4}$,
\begin{equation}
S_{\mathrm{int}} \supset
\frac{1}{M^4} \int d^4x \sqrt{-g} \,
T^{(+)}_{\mu\nu} T_{(-)}^{\mu\nu}.
\end{equation}

This correction induces a small modulation of the tensor power spectrum,
\begin{equation}
\Omega_{\mathrm{GW}}(k)
=
\Omega_{\mathrm{GW}}^{(0)}(k)
\left[ 1 + \alpha_{\mathrm{IR}} \log(k/k_0) \right],
\end{equation}
and a quadratic suppression,
\begin{equation}
\Omega_{\mathrm{GW}} \propto \alpha_{\mathrm{IR}}^2.
\end{equation}

These signatures are unique to RDCC and arise from the global coherence of
tensor modes.
% ========================================================
% 4 CONCLUSION
% ========================================================

\section{Conclusion}

This Companion develops the corrected tensor sector of RDCC in its v49.1 form.
Tensor modes are global geometric degrees of freedom that remain coherent across
both sectors of the CPT-symmetric bipartite Hilbert space. Their evolution is
governed by the global structure of the metric and the ultraviolet-suppressed
dimension-8 inter-sector operator.

The resulting quadratic suppression and logarithmic modulation of the tensor
spectrum form the basis of RDCC's gravitational-wave phenomenology. These
results integrate the relaxation sector (Companion III), the ultraviolet sector
(Companion X), and the tensor–relaxation coupling developed in Companion XI.
Observational predictions are presented in Companion XIV.

\end{document}
